\chapter{Esercizi: Complessità computazionale}
\label{chap:esercizi-complessita}

\begin{exercise}
\label{ex:fi-esercizi-complessita-tipi-analisi}
Quali sono i due tipi di analisi di un algoritmo? In che cosa
si differenziano?
\end{exercise}

Si distinguono l'\emph{analisi sperimentale}, che misura le prestazioni
di un'implementazione su ingressi concreti, e l'\emph{analisi teorica},
che determina il costo delle operazioni in un modello astratto.
La prima dipende dall'ambiente di esecuzione; la seconda permette
di studiare la crescita delle risorse richieste senza effettuare
misurazioni (\cref{sec:fi-complessita-sperimentale,sec:fi-complessita-modello}).

\begin{exercise}
\label{ex:fi-esercizi-complessita-sperimentale-implementazione}
Perché diciamo che l'analisi sperimentale misura le proprietà di
un'\emph{implementazione} e non dell'\emph{algoritmo} in quanto tale?
Elenca almeno tre fattori che influenzano i tempi misurati
sperimentalmente.
\end{exercise}

Le misure riguardano l'esecuzione di un programma concreto e non
il procedimento astratto. I tempi dipendono, tra l'altro,
dall'hardware, dal linguaggio e dall'interprete o compilatore,
dai valori dell'ingresso e dalle condizioni di esecuzione.
Implementazioni diverse dello stesso algoritmo possono quindi
produrre tempi differenti (\cref{sec:fi-complessita-sperimentale}).

\begin{exercise}
\label{ex:fi-esercizi-complessita-raddoppio}
In un esperimento di scalabilità, perché conviene raddoppiare
la dimensione dell'ingresso a ogni prova invece di aumentarla
di una quantità fissa?
\end{exercise}

Il raddoppio permette di osservare la crescita su un intervallo
ampio di dimensioni con poche misurazioni. Se il tempo segue
il modello \(t(n)=an^b\), il rapporto \(t(2n)/t(n)=2^b\) è costante
e consente di stimare l'esponente \(b\). Un incremento fisso
fornisce invece rapporti che dipendono da \(n\)
(\cref{sec:fi-complessita-sperimentale}).

\begin{exercise}
\label{ex:fi-esercizi-complessita-stima-esponente}
In un esperimento i tempi di esecuzione misurati sono
\(T(n)=3\,\mathrm{s}\) e \(T(2n)=24\,\mathrm{s}\). Stima l'esponente
\(b\) nella relazione \(T(n)\approx a\,n^b\). Che cosa puoi dire
della complessità dell'algoritmo?
\end{exercise}

Si ricava
\[
  \frac{T(2n)}{T(n)}=\frac{24}{3}=8=2^b,
  \qquad b=\log_2 8=3.
\]
La misura suggerisce quindi una crescita \emph{cubica},
\(T(n)\propto n^3\). Due sole misurazioni non dimostrano però
che la complessità asintotica sia \(\Theta(n^3)\): per verificarne
la plausibilità occorrono più prove, mentre una dimostrazione
richiede l'analisi teorica
(\cref{sec:fi-complessita-sperimentale,sec:fi-complessita-modello}).

\begin{exercise}
\label{ex:fi-esercizi-complessita-passo-base}
Che cosa si intende per ``passo base'' (o operazione primitiva)?
Fai due esempi.
\end{exercise}

Un \emph{passo base} è un'operazione cui si attribuisce costo
costante nel modello adottato. Sono esempi l'assegnamento
di un valore a una variabile e il confronto fra due interi
di dimensione limitata. La scelta dei passi base deve essere
coerente con le ipotesi del modello
(\cref{sec:fi-complessita-modello}).

\begin{exercise}
\label{ex:fi-esercizi-complessita-chiamata-funzione}
Perché l'invocazione di una funzione in generale non può essere
considerata un'operazione primitiva?
\end{exercise}

Il costo di una chiamata comprende le operazioni svolte dalla
funzione, che possono dipendere dalla dimensione dell'ingresso.
Per esempio, \texttt{sorted(lista)} ordina gli elementi della lista:
la sua invocazione non ha costo costante rispetto alla lunghezza
dell'ingresso (\cref{sec:fi-complessita-modello}).

\begin{exercise}
\label{ex:fi-esercizi-complessita-passi-contesto}
I passi considerati ``base'' sono sempre gli stessi, oppure
possono cambiare a seconda del contesto? Motiva la risposta.
\end{exercise}

Dipendono dal \emph{modello di costo} e dal livello di dettaglio
richiesto. Il modello a costo unitario considera costanti alcune
operazioni elementari su operandi di dimensione limitata; se
invece si elaborano interi arbitrariamente grandi, il costo
delle operazioni aritmetiche può dipendere dal numero di bit.
Le ipotesi devono essere dichiarate prima dell'analisi
(\cref{sec:fi-complessita-modello,sec:fi-complessita-fibonacci}).

\begin{exercise}
\label{ex:fi-esercizi-complessita-definizione-o}
Scrivi la definizione formale di \(f(n)\in O(g(n))\).
Qual è il ruolo delle costanti \(c\) e \(n_0\)?
\end{exercise}

Per funzioni \(f,g\colon\N\to[0,+\infty)\), con \(g(n)>0\)
per \(n\) sufficientemente grande,
\[
  f\in O(g)
  \iff
  \exists c>0\ \exists n_0\in\N:\
  f(n)\le c\,g(n)\quad\forall n\ge n_0.
\]
La costante \(c\) stabilisce il fattore moltiplicativo del
limite superiore, mentre \(n_0\) indica da quale dimensione
tale limite deve valere
(\cref{def:fi-notazioni-asintotiche}).

\begin{exercise}
\label{ex:fi-esercizi-complessita-confronto-notazioni}
Spiega a parole la differenza tra \(O\), \(\Omega\) e \(\Theta\).
\end{exercise}

Le notazioni \(O\) e \(\Omega\) esprimono rispettivamente un
limite superiore e uno inferiore asintotico, a meno di
costanti moltiplicative positive. La notazione \(\Theta\) esprime
entrambi i limiti e quindi una stima asintotica dello stesso
ordine di crescita
(\cref{def:fi-notazioni-asintotiche,fig:fi-notazioni-asintotiche}).

\begin{exercise}
\label{ex:fi-esercizi-complessita-vero-falso}
Stabilisci se le seguenti affermazioni sono vere o false,
motivando la risposta:
\[
  5n^2+100n\in\Theta(n^2),\qquad
  n\in O(n^2),
\]
\[
  n^2\in\Theta(n),\qquad
  2^n+n^{10}\in O(2^n).
\]
\end{exercise}

La prima è \emph{vera}: per \(n\ge1\),
\(5n^2\le5n^2+100n\le105n^2\).
La seconda è \emph{vera}, poiché \(n\le n^2\) per \(n\ge1\).
La terza è \emph{falsa}: il rapporto \(n^2/n=n\) non è
superiormente limitato da una costante.
La quarta è \emph{vera}: la funzione esponenziale \(2^n\)
cresce più rapidamente di \(n^{10}\), quindi
\(n^{10}\in O(2^n)\) e \(2^n+n^{10}\in O(2^n)\)
(\cref{def:fi-notazioni-asintotiche,sec:fi-complessita-ordini}).

\begin{exercise}
\label{ex:fi-esercizi-complessita-omega-cubico}
È vero che \(n\in\Omega(n^2)\)? E che \(n^3\in O(n^3)\)?
Giustifica.
\end{exercise}

La prima affermazione è \emph{falsa}: per \(n>0\), il rapporto
\(n/n^2=1/n\) tende a zero, quindi non esiste una costante
positiva \(c\) tale che \(n\ge cn^2\) per tutti gli \(n\)
sufficientemente grandi. La seconda è \emph{vera}: basta
scegliere \(c=1\) nella definizione di \(O\)
(\cref{def:fi-notazioni-asintotiche}).

\begin{exercise}
\label{ex:fi-esercizi-complessita-regole-costo}
Enuncia le regole per derivare la complessità di un blocco
di codice sequenziale, di un \lstinline|if| e di un ciclo.
\end{exercise}

In una sequenza si sommano i costi dei blocchi. In un
\lstinline|if| si somma il costo di valutazione della condizione
al costo del ramo eseguito; per un limite superiore si può
usare quello del ramo più oneroso, aggiungendo sempre il costo
della condizione. Per un ciclo si sommano i costi delle singole
iterazioni, che possono essere diversi. Se ogni iterazione
costa \(\Theta(1)\), il costo del corpo è proporzionale al
numero di iterazioni, salvo gli eventuali costi aggiuntivi
della struttura di controllo
(\cref{sec:fi-complessita-modello}).

\clearpage

\begin{exercise}
\label{ex:fi-esercizi-complessita-costante}
Perché, una volta stabilito che un blocco di codice ha costo
costante, il valore preciso di quella costante è irrilevante
ai fini della classe di complessità?
\end{exercise}

Le notazioni asintotiche trascurano i fattori moltiplicativi
costanti: qualunque costo positivo indipendente da \(n\)
appartiene a \(\Theta(1)\). Il valore della costante resta però
rilevante per il tempo effettivo su ingressi concreti
(\cref{sec:fi-complessita-asintotica}).

\begin{exercise}
\label{ex:fi-esercizi-complessita-cicli-quadrati}
Considera il seguente frammento di codice Python. Qual è
la sua complessità in funzione di \(n\)? Deriva il numero
di operazioni eseguite dall'istruzione più interna.
\end{exercise}

\begin{unipdcode}[language=Python]
total = 0
for i in range(n):
    for j in range(n):
        total = total + 1
\end{unipdcode}

Per ogni valore di \(i\), il ciclo interno esegue \(n\) volte
l'aggiornamento di \texttt{total}. Il ciclo esterno ha \(n\)
iterazioni, quindi l'istruzione più interna viene eseguita
\[
  \sum_{i=0}^{n-1}n=n^2
\]
volte. Nel modello a costo unitario il tempo è
\(\Theta(n^2)\) (\cref{sec:fi-complessita-modello}).

\begin{exercise}
\label{ex:fi-esercizi-complessita-ciclo-triangolare}
Considera ora il seguente frammento, in cui il ciclo interno
dipende da \lstinline|i|. Quante volte viene eseguita
l'istruzione più interna, in totale? Qual è la complessità?
\end{exercise}

\begin{unipdcode}[language=Python]
total = 0
for i in range(n):
    for j in range(i):
        total = total + 1
\end{unipdcode}

Quando \(i=0\) il ciclo interno non viene eseguito; per un
indice generico \(i\), compie \(i\) iterazioni. Il numero
totale di aggiornamenti è
\[
  \sum_{i=0}^{n-1}i=\frac{n(n-1)}2.
\]
Il termine dominante è quadratico, quindi la complessità
è \(\Theta(n^2)\) (\cref{sec:fi-complessita-modello}).

\begin{exercise}
\label{ex:fi-esercizi-complessita-reverse-quadratico}
Perché il primo algoritmo \texttt{reverse}, presentato nella
\cref{sec:fi-complessita-reverse-quadratico}, ha complessità
\(\Theta(n^2)\)? Qual è il ``lavoro nascosto'' che lo rende
quadratico?
\end{exercise}

Ogni nuovo elemento viene inserito all'inizio della lista di
uscita, spostando verso destra gli \(i\) elementi già inseriti.
Il numero complessivo degli spostamenti è
\[
  \sum_{i=0}^{n-1}i=\frac{n(n-1)}2.
\]
Nel modello di conteggio degli appunti si ottiene
\(T(n)=\frac32n^2+\frac52n+3\). Il costo quadratico dipende
quindi dagli \emph{spostamenti ripetuti}, non dal semplice
numero di iterazioni del ciclo esterno
(\cref{sec:fi-complessita-reverse-quadratico}).

\begin{exercise}
\label{ex:fi-esercizi-complessita-reverse-lineare}
Qual è l'idea che permette al secondo algoritmo
\texttt{reverse}, presentato nella
\cref{sec:fi-complessita-reverse-lineare}, di raggiungere
una complessità \(\Theta(n)\)?
\end{exercise}

Ogni elemento viene copiato direttamente nella sua
posizione definitiva della lista invertita: all'indice \(i\)
si scrive l'elemento di indice \(n-1-i\) dell'ingresso.
Si evitano così gli spostamenti degli elementi già copiati.
Il costo è \(T(n)=5n+2\) nel modello degli appunti
(\cref{sec:fi-complessita-reverse-lineare}).

\begin{exercise}
\label{ex:fi-esercizi-complessita-rapporto-reverse}
Supponi che per \(n=32000\) l'algoritmo lineare risulti circa
12\,800 volte più veloce di quello quadratico. Questo rapporto rimane costante
al crescere di \(n\), oppure cambia? Perché?
\end{exercise}

Il rapporto non rimane costante. I costi dei due algoritmi
sono rispettivamente \(\Theta(n^2)\) e \(\Theta(n)\), quindi,
a parità di modello, il rapporto fra il costo quadratico e
quello lineare cresce come \(\Theta(n)\). Raddoppiando \(n\),
il rapporto tende approssimativamente a raddoppiare.
Il valore 12\,800 è un dato del quesito e non si ricava
dalle sole classi asintotiche: i tempi effettivi dipendono
anche dall'implementazione e dall'ambiente
(\cref{sec:fi-complessita-reverse-quadratico,sec:fi-complessita-reverse-lineare,sec:fi-complessita-sperimentale}).

\begin{exercise}
\label{ex:fi-esercizi-complessita-computer-veloce}
``Con un computer abbastanza veloce, l'algoritmo quadratico
diventa preferibile a quello lineare.'' Questa affermazione
è corretta? Discuti.
\end{exercise}

No, non in generale. Un computer più veloce può diminuire
i tempi di entrambi gli algoritmi, ma non modifica il
loro ordine di crescita. Per ingressi sufficientemente
grandi, la crescita quadratica rende il primo algoritmo
asintoticamente meno efficiente di quello lineare.
Per dimensioni concrete possono contare le costanti e
i dettagli dell'implementazione, che non cambiano però
il confronto asintotico
(\cref{sec:fi-complessita-asintotica,sec:fi-complessita-ordini,sec:fi-complessita-reverse-lineare}).
